\documentclass[aspectratio=169,10pt]{beamer}

\usetheme{default}
\setbeamertemplate{navigation symbols}{}
\setbeamersize{text margin left=0.65cm,text margin right=0.65cm}

\usepackage[T1]{fontenc}
\usepackage{lmodern}
\usepackage{graphicx}
\usepackage{booktabs}
\usepackage{array}
\usepackage{amsmath}
\usepackage{url}

\setbeamercolor{normal text}{fg=black,bg=white}
\setbeamercolor{frametitle}{fg=black,bg=white}
\setbeamerfont{frametitle}{series=\bfseries,size=\Large}
\setbeamertemplate{itemize item}{--}
\setbeamertemplate{itemize subitem}{--}
\setbeamertemplate{footline}{%
  \hfill{\scriptsize Compact magnetometer update \insertframenumber/\inserttotalframenumber}\hspace{0.6em}\vspace{0.35em}
}

\newcommand{\assetdir}{assets/project_update_20260514}
\newcommand{\smallitemsep}{\setlength{\itemsep}{0.45em}}

\title{Compact Magnetometer Project Update}
\subtitle{Updated lock-in and live B-field data}
\author{Chalach Kasemtantikul}
\date{May 14, 2026}

\begin{document}

\begin{frame}{Compact magnetometer project update}
\begin{columns}[T,totalwidth=\textwidth]
\begin{column}{0.56\textwidth}
\begin{itemize}\smallitemsep
  \item The lock-in workflow is still the same:
  \[
  \begin{aligned}
  \text{ODMR sweep} &\rightarrow \text{16 parked frequencies} \\
  &\rightarrow \Delta f(t) \rightarrow \Delta B(t).
  \end{aligned}
  \]
  \item May 14 update: spike rejection is now active in live lock-in and live B-field acquisition.
  \item Latest parking run: 8 kept transitions, 16 parked frequencies, no rejected fitted centers.
  \item Live lock-in test: 266 batches in 21.3 s, about 12.5 Hz update rate.
  \item Offline B reconstruction: 266/266 samples rank-ok, condition number 9.74.
  \item Live B-field run: 387 batches in 60.0 s, about 6.4 Hz update rate.
\end{itemize}
\end{column}
\begin{column}{0.40\textwidth}
\centering
\begin{tabular}{@{}ll@{}}
\toprule
Item & Current result \\
\midrule
Transitions & 8 \\
Parked points & 16 \\
Single batch & 113 ms \\
Live lock-in & 12.5 Hz \\
Live B-field & 6.4 Hz \\
B matrix cond. & 9.74 \\
\bottomrule
\end{tabular}

\vspace{1.2em}
\footnotesize
Latest live lock-in spike replacements:

\vspace{0.2em}
\begin{tabular}{@{}lr@{}}
\toprule
Signal cells & 48 \\
Reference cells & 24 \\
\bottomrule
\end{tabular}
\end{column}
\end{columns}
\end{frame}

\begin{frame}{Workflow 1: choose 16 parked frequencies}
\begin{columns}[T,totalwidth=\textwidth]
\begin{column}{0.70\textwidth}
\centering
\includegraphics[width=\linewidth,height=0.72\textheight,keepaspectratio]{\assetdir/workflow1_parked_frequencies.png}
\end{column}
\begin{column}{0.27\textwidth}
\begin{itemize}\smallitemsep
  \item Latest ODMR sweep: May 14 12:09 run.
  \item The new parking algorithm kept all 8 detected transitions.
  \item Two parked frequencies are placed on each transition.
  \item The estimated bias field from the ODMR fit is approximately $(2.12,-0.59,-3.42)$ mT.
\end{itemize}
\end{column}
\end{columns}
\end{frame}

\begin{frame}{Workflow 2: acquire all parked points}
\begin{columns}[T,totalwidth=\textwidth]
\begin{column}{0.72\textwidth}
\centering
\includegraphics[width=\linewidth,height=0.72\textheight,keepaspectratio]{\assetdir/workflow2_parked_overlay.png}
\end{column}
\begin{column}{0.25\textwidth}
\footnotesize
\begin{itemize}\setlength{\itemsep}{0.2em}
  \item One FPGA program measures all 16 parked frequencies.
  \item The one-batch reference is used as the baseline for live despiking.
  \item Latest one-batch acquisition took 113 ms.
  \item The predicted FPGA work per live batch is about 69 ms for 10 reps.
\end{itemize}
\end{column}
\end{columns}
\end{frame}

\begin{frame}{Workflow 3: convert lock-in signal to frequency shift}
\begin{columns}[T,totalwidth=\textwidth]
\begin{column}{0.49\textwidth}
\footnotesize
For each transition, use the two flanks:
\[
\Delta f_k(t)=
\frac{\left[S_+(t)-S_-(t)\right]-\left[S_{+,0}-S_{-,0}\right]}
{m_- - m_+}.
\]
\vspace{0.2em}
\begin{itemize}\setlength{\itemsep}{0.25em}
  \item $S_-$ and $S_+$ are the live signals at the lower and upper parked points.
  \item $S_{-,0}$ and $S_{+,0}$ come from the reference ODMR calibration.
  \item Latest live run: 266 batches over 21.3 s.
  \item Per-block frequency noise: 0.09 to 0.45 MHz standard deviation.
  \item Projected-field noise: 3.2 to 16.6 $\mu$T standard deviation.
\end{itemize}
\end{column}
\begin{column}{0.48\textwidth}
\centering
\includegraphics[width=\linewidth,height=0.70\textheight,keepaspectratio]{\assetdir/workflow3_delta_f.png}
\end{column}
\end{columns}
\end{frame}

\begin{frame}{Workflow 4: inferred new peak frequencies}
\begin{columns}[T,totalwidth=\textwidth]
\begin{column}{0.66\textwidth}
\centering
\includegraphics[width=\linewidth,height=0.78\textheight,keepaspectratio]{\assetdir/workflow4_inferred_peaks.png}
\end{column}
\begin{column}{0.31\textwidth}
\begin{itemize}\smallitemsep
  \item This is the inferred new peak frequency for each of the 8 transitions.
  \item Solid traces show $f_{\mathrm{new}}(t)=f_0+\Delta f(t)$.
  \item The May 14 live run uses the causal Hampel despiker before the frequency conversion.
  \item This view is still useful for spotting remaining jumps before B reconstruction.
\end{itemize}
\end{column}
\end{columns}
\end{frame}

\begin{frame}{Workflow 5: reconstruct magnetic field}
\begin{columns}[T,totalwidth=\textwidth]
\begin{column}{0.76\textwidth}
\centering
\includegraphics[width=\linewidth,height=0.49\textheight,keepaspectratio]{\assetdir/workflow5_b_components.png}

\vspace{0.15em}
\includegraphics[width=\linewidth,height=0.23\textheight,keepaspectratio]{\assetdir/workflow5_b_magnitude.png}
\end{column}
\begin{column}{0.21\textwidth}
\scriptsize
\begin{itemize}\setlength{\itemsep}{0.18em}
  \item Offline reconstruction of the May 14 no-magnet run.
  \item 266 samples, 266 rank-ok.
  \item B matrix condition number: 9.74.
  \item Reconstructed B statistics:
\end{itemize}
\vspace{0.1em}
\centering
\tiny
\begin{tabular}{@{}lrr@{}}
\toprule
Comp. & Mean & Std. \\
\midrule
$\Delta B_x$ & 24.51 & 4.49 \\
$\Delta B_y$ & -11.38 & 6.42 \\
$\Delta B_z$ & 29.91 & 6.53 \\
$|\Delta B|$ & 40.91 & 7.41 \\
\bottomrule
\end{tabular}

\vspace{0.15em}
\tiny Units: $\mu$T
\end{column}
\end{columns}
\end{frame}

\begin{frame}{New data: small magnet}
\begin{columns}[T,totalwidth=\textwidth]
\begin{column}{0.70\textwidth}
\centering
\includegraphics[width=\linewidth,height=0.80\textheight,keepaspectratio]{\assetdir/small_b_reconstruction.png}
\end{column}
\begin{column}{0.27\textwidth}
\footnotesize
\begin{itemize}\setlength{\itemsep}{0.25em}
  \item Source: May 14 small-magnet run.
  \item 318 samples over 25.9 s.
  \item Live lock-in rate: 12.3 Hz.
  \item B matrix condition number: 9.74.
\end{itemize}
\vspace{0.3em}
\centering
\scriptsize
\begin{tabular}{@{}lrr@{}}
\toprule
Comp. & Mean & Std. \\
\midrule
$\Delta B_x$ & 25.19 & 7.23 \\
$\Delta B_y$ & -10.93 & 9.91 \\
$\Delta B_z$ & 29.25 & 8.11 \\
$|\Delta B|$ & 41.68 & 9.38 \\
\bottomrule
\end{tabular}

\vspace{0.15em}
\scriptsize Units: $\mu$T
\end{column}
\end{columns}
\end{frame}

\begin{frame}{New data: neodymium big magnet}
\begin{columns}[T,totalwidth=\textwidth]
\begin{column}{0.70\textwidth}
\centering
\includegraphics[width=\linewidth,height=0.80\textheight,keepaspectratio]{\assetdir/neodymium_big_b_reconstruction.png}
\end{column}
\begin{column}{0.27\textwidth}
\footnotesize
\begin{itemize}\setlength{\itemsep}{0.25em}
  \item Source: May 14 neodymium run.
  \item 344 samples over 26.3 s.
  \item Live lock-in rate: 13.1 Hz.
  \item This run shows larger transient spikes than the small-magnet run.
\end{itemize}
\vspace{0.3em}
\centering
\scriptsize
\begin{tabular}{@{}lrr@{}}
\toprule
Comp. & Mean & Std. \\
\midrule
$\Delta B_x$ & 26.08 & 18.08 \\
$\Delta B_y$ & -10.27 & 27.22 \\
$\Delta B_z$ & 31.76 & 24.34 \\
$|\Delta B|$ & 49.37 & 31.87 \\
\bottomrule
\end{tabular}

\vspace{0.15em}
\scriptsize Units: $\mu$T
\end{column}
\end{columns}
\end{frame}

\begin{frame}{New data: offline comparison}
\centering
\includegraphics[width=0.92\textwidth,height=0.58\textheight,keepaspectratio]{\assetdir/offline_case_b_stats_bar.png}

\vspace{0.4em}
\footnotesize
\begin{itemize}\setlength{\itemsep}{0.2em}
  \item The comparison includes small magnet, neodymium big magnet, and no-magnet baseline.
  \item The neodymium run has the largest standard deviation because of transient jumps.
  \item Reporting both mean and standard deviation helps separate field offset from stability.
\end{itemize}
\end{frame}

\begin{frame}{Workflow 6: live B-field measurement}
\begin{columns}[T,totalwidth=\textwidth]
\begin{column}{0.74\textwidth}
\centering
\includegraphics[width=\linewidth,height=0.82\textheight,keepaspectratio]{\assetdir/workflow6_live_b_latest.png}
\end{column}
\begin{column}{0.23\textwidth}
\scriptsize
\begin{itemize}\setlength{\itemsep}{0.15em}
  \item Source: May 14 live-B run.
  \item 387 batches in 60.0 s.
  \item Update rate: about 6.4 Hz.
  \item 387/387 B reconstructions rank-ok.
  \item Spike replacements: 415 signal, 404 reference.
\end{itemize}
\vspace{0.2em}
\centering
\tiny
\begin{tabular}{@{}lrr@{}}
\toprule
Comp. & Mean & Std. \\
\midrule
$\Delta B_x$ & 23.72 & 12.76 \\
$\Delta B_y$ & 0.77 & 29.79 \\
$\Delta B_z$ & 31.60 & 17.14 \\
$|\Delta B|$ & 49.28 & 21.78 \\
\bottomrule
\end{tabular}

\vspace{0.15em}
\scriptsize Units: $\mu$T
\end{column}
\end{columns}
\end{frame}

\begin{frame}{Live B magnet comparison}
\begin{columns}[T,totalwidth=\textwidth]
\begin{column}{0.68\textwidth}
\centering
\includegraphics[width=\linewidth,height=0.78\textheight,keepaspectratio]{\assetdir/live_b_total_comparison.png}
\end{column}
\begin{column}{0.29\textwidth}
\footnotesize
\begin{tabular}{@{}lrr@{}}
\toprule
Run & Mean & Std. \\
\midrule
3-stack & 50.83 & 46.24 \\
Neodymium & 49.28 & 21.78 \\
\bottomrule
\end{tabular}

\vspace{0.15em}
\scriptsize Units: $\mu$T for $|\Delta B|$

\vspace{0.8em}
\begin{itemize}\setlength{\itemsep}{0.25em}
  \item Both live-B test runs operated at about 6.5 Hz.
  \item The 3-stack run has the larger variation.
  \item Next step is to compare these live-B outputs against a known field geometry.
\end{itemize}
\end{column}
\end{columns}
\end{frame}

\begin{frame}{Next steps}
\begin{columns}[T,totalwidth=\textwidth]
\begin{column}{0.48\textwidth}
\textbf{Validate B reconstruction}
\begin{itemize}\smallitemsep
  \item Apply a known test field and compare recovered direction and magnitude.
  \item Repeat no-magnet, small-magnet, and neodymium measurements with fixed geometry.
  \item Track residual norm, matrix rank, and condition number for every batch.
  \item Check whether the spike rejection removes artifacts without hiding real field changes.
\end{itemize}
\end{column}
\begin{column}{0.48\textwidth}
\textbf{Improve live operation}
\begin{itemize}\smallitemsep
  \item Move live B closer to the 12.5 Hz raw live lock-in rate.
  \item Decouple acquisition, reconstruction, and plotting.
  \item Continue dual-channel measurement integration.
  \item Continue hardware optimization and patent manuscript documentation.
\end{itemize}
\end{column}
\end{columns}
\end{frame}

\end{document}
